```latex
\documentclass[11pt]{article}
\usepackage{amsmath,amssymb,amsthm,geometry,mathtools}
\usepackage{hyperref}
\geometry{margin=1.1in}

\newtheorem{theorem}{Theorem}[section]
\newtheorem{lemma}[theorem]{Lemma}
\newtheorem{proposition}[theorem]{Proposition}
\newtheorem{corollary}[theorem]{Corollary}
\newtheorem{remark}[theorem]{Remark}
\newtheorem{definition}[theorem]{Definition}
\newtheorem{assumption}[theorem]{Assumption}

\newcommand{\R}{\mathbb{R}}
\newcommand{\E}{\mathbb{E}}
\newcommand{\W}{\mathbb{W}}
\newcommand{\1}{\mathbf{1}}
\newcommand{\norm}[1]{\lVert #1\rVert}
\newcommand{\inner}[2]{\langle #1,#2\rangle}
\newcommand{\ddd}{\mathrm{d}}
\newcommand{\Lto}[2]{#1\to #2}
\newcommand{\Gf}{\mathcal{G}}
\newcommand{\F}{\mathcal{F}}
\DeclareMathOperator{\divg}{div}
\DeclareMathOperator{\tr}{tr}
\DeclareMathOperator{\diag}{diag}

\title{Minimax Estimation in the Statistical Darcy Problem on the Unit Ball:\\
The Role of Radial Transport and Anisotropic Ill--Posedness}
\author{}
\date{}

\begin{document}
\maketitle

\begin{abstract}
We consider the nonlinear statistical inverse problem of recovering a conductivity
$f\in C^\alpha(D)$, $\alpha>1$, on the unit ball $D\subset\R^d$, $d\ge 2$, from a Gaussian
white--noise observation of the solution $u_f$ of the Dirichlet problem
$\divg(f\nabla u)=1$ in $D$, $u|_{\partial D}=0$, over a H\"older ball
$\F(\alpha,r)=\{f:\norm{f-\1}_{C^\alpha}\le r\}$ of small fixed radius $r\in(0,1)$.
We prove that the minimax risk in the $L^2$--mean squared error satisfies, up to
logarithmic factors,
\[
\inf_{\hat f_n}\sup_{f_0\in\F(\alpha,r)}\E_{f_0}\norm{\hat f_n-f_0}^2_{L^2(D)}
\;\asymp\;
r^{\frac{2(d+3)}{2\alpha+d+3}}\; n^{-\frac{2\alpha}{2\alpha+d+3}},
\qquad n\to\infty .
\]
In particular the minimax \emph{rate} is $n^{-2\alpha/(2\alpha+d+3)}$ up to
polylogarithmic factors. The exponent $d+3$ (as opposed to $d+4$, which would
correspond to an isotropically $2$--ill--posed problem, or to the rates
$n^{-2(\alpha-1)/(2\alpha+d+2)}$ obtained in the boundary--constrained setting
of Nickl--van de Geer and subsequent work) reflects an \emph{anisotropic} degree of
ill--posedness: at the linearization around $f\equiv 1$ the forward map smooths by two
derivatives in the $d-1$ angular (spherical) directions but by only \emph{one} derivative
along the radial transport generated by the unique critical point of $u_f$ near the origin.
Both the lower bound (an Assouad construction based on angular wave packets) and the
upper bound (a transport--inverted spectral estimator exploiting an explicit left inverse
of the linearized forward map) are given.
\end{abstract}

\tableofcontents

\section{Introduction and main result}\label{sec:intro}

\subsection{The statistical model}

Let $D=\{x\in\R^d:|x|<1\}$, $d\ge 2$, with $d\ge 2$ essential throughout. For a
conductivity $f>0$ on $D$ let $u_f$ denote the unique weak solution of Darcy's problem
\begin{equation}\label{eq:pde}
\divg(f\nabla u)=1\ \text{ in } D,\qquad u|_{\partial D}=0 .
\end{equation}
For smoothness parameter $\alpha>1$ and $r\in(0,1)$ consider
\[
\F(\alpha,r)=\big\{h\in C^\alpha(\bar D):\ \norm{h-\1}_{C^\alpha(\bar D)}\le r\big\},
\]
where $C^\alpha$ denotes the usual H\"older--Zygmund scale. Note that we do
\emph{not} impose any boundary constraint on $f$ beyond $f\in \F(\alpha,r)$; this
will be seen to be consequential. We observe, in the Gaussian white--noise model on
$L^2(D)$,
\begin{equation}\label{eq:wn}
Y_n=u_f+\frac{1}{\sqrt n}\,\W,\qquad\qquad
\inner{Y_n}{\varphi}=\inner{u_f}{\varphi}+\frac{1}{\sqrt n}\int_0^1\!\inner{\varphi}{\ddd W_t}\,\ddd t ,
\end{equation}
i.e.\ $Y_n$ is the Gaussian random element in $L^2(D)$ with mean $u_f$ and covariance
$n^{-1}\mathrm{Id}$. We study the minimax risk
\begin{equation}\label{eq:minimax}
\mathcal R_n(\alpha,r)\ :=\ \inf_{\hat f_n}\ \sup_{f_0\in\F(\alpha,r)}
\E_{f_0}\Big[\norm{\hat f_n-f_0}^2_{L^2(D)}\Big],
\end{equation}
where the infimum extends over all measurable functions of the data \eqref{eq:wn}.

\begin{theorem}[Main theorem]\label{thm:main}
Let $d\ge 2$, $\alpha>1$, and $r\in(0,r_0]$ for some $r_0=r_0(\alpha,d)\in(0,1/2)$
sufficiently small. Then there exist constants $0<c\le C<\infty$ depending only on
$(\alpha,d,r)$, and integers $c_1,c_2$ depending only on $(\alpha,d)$, such that for all
$n\ge 2$,
\[
c\,r^{\frac{2(d+3)}{2\alpha+d+3}}\,n^{-\frac{2\alpha}{2\alpha+d+3}}
\ \le\ \mathcal R_n(\alpha,r)\ \le\
C\,(\log n)^{c_2}\,r^{\frac{2(d+3)}{2\alpha+d+3}}\,n^{-\frac{2\alpha}{2\alpha+d+3}} .
\]
Consequently, for every fixed (small) $r$, the minimax rate for the $L^2$--mean squared
error, up to logarithmic factors, is
\begin{equation}\label{eq:rate}
\boxed{\ \mathcal R_n(\alpha,r)\ \asymp\ n^{-\frac{2\alpha}{2\alpha+d+3}}\ }
\end{equation}
(with the $r$--dependence $r^{2(d+3)/(2\alpha+d+3)}$ as displayed above).
\end{theorem}

\subsection{Discussion and relation to the literature}

Statistical inference for PDE--constrained inverse problems of the type
\eqref{eq:pde}--\eqref{eq:wn}, in regression designs, has been developed in a series of
papers: Tikhonov--type penalised least squares (MAP/Bayesian estimators) were analysed
for divergence--form and Schr\"odinger equations in
\href{https://export.arxiv.org/pdf/1809.08818v3.pdf}{arXiv:1809.08818} (Nickl--van de Geer),
MAP estimators for random design in
\href{https://arxiv.org/html/2409.03417}{arXiv:2409.03417}, Bayesian nonparametric
procedures with Gaussian process priors in
\href{https://google.iopscience.iop.org/article/10.1088/1361-6420/ab7d2a}{Inverse Problems 36 (2020)},
and computationally efficient (polynomial time) estimators achieving the same rates in
\href{https://arxiv.org/html/2601.09007}{arXiv:2601.09007}. In that literature the
parameter class \emph{fixes the boundary values} (e.g.\ $f|_{\partial D}\equiv 1$ together
with vanishing normal derivatives up to order $\lfloor\alpha-1\rfloor$), which is needed
for the identifiability/stability arguments based on the $2$--ill--posed estimate
$\norm{f_1-f_2}_{L^2}\lesssim\norm{u_{f_1}-u_{f_2}}_{H^2}$, and the resulting recovery
rate for $f$ is $N^{-2(\alpha-1)/(2\alpha+d+2)}$, which is \emph{not} claimed to be minimax
for $f$.

The class $\F(\alpha,r)$ of Theorem~\ref{thm:main} deliberately does \emph{not} fix
boundary values. The analysis then has to face the true local geometry of the forward
map, and it turns out that the problem is \emph{anisotropically} mildly ill--posed:
\begin{itemize}
\item in the $d-1$ \emph{angular} (spherical) directions, the linearised forward map
attenuates a frequency of magnitude $m$ like $m^{-2}$ (degree $2$ ill--posedness);
\item along the \emph{radial} direction, the transport structure generated by the
equation $x\cdot\nabla h+d\,h=\text{source}$ (cf.\ Lemma~\ref{lem:transport}) means the
attenuation is only of first order.
\end{itemize}
The hardest subproblems are therefore ``angular wave packets, constant along radii'';
there are of order $L^{d-1}$ of them at total frequency $L$, each carrying inverse
attenuation $L^4$. This yields a variance exponent $L^{d+3}$ for any spectral cut--off
level $L$, and hence, balancing against the H\"older bias $r^2L^{-2\alpha}$, the rate
\eqref{eq:rate}. We emphasise:
\begin{equation*}
\underbrace{2\alpha+d+3}_{\text{this paper}}
\quad<\quad
2\alpha+d+4\ \ (\text{isotropically }2\text{--ill--posed}),
\end{equation*}
so exploiting the radial transport \emph{strictly improves} the rate; and the
unconstrained problem is \emph{harder} than the boundary--constrained one treated in
the above references, in the sense that our lower bound applies a fortiori to larger
classes and our upper bound gives a faster decay than the constrained--class upper
bounds when $\alpha>d+3$ (there is no contradiction: the constrained classes are
strictly smaller, and the published rates for them are upper bounds).

Throughout, $C,c$ denote constants possibly depending on $(\alpha,d,r)$ but never on $n$;
$\lesssim,\gtrsim,\asymp$ are understood up to such constants. For the elliptic
regularity background we cite \cite{GT}; for H\"older--Zygmund and Besov spaces and
spectral approximation numbers \cite{Triebel}; for Bessel function asymptotics
\cite{Watson,Olver}; for Assouad's lemma and general minimax lower bound technique
\cite{Tsybakov}.

\section{Preliminaries: the forward map}\label{sec:forward}

For $r\le 1/2$ every $f\in\F(\alpha,r)$ satisfies $f\ge 1-r\ge 1/2$, so
\eqref{eq:pde} is uniformly elliptic. Define the forward map
\[
\Gf:\ \F(\alpha,r)\to L^2(D),\qquad \Gf(f)=u_f .
\]
For $f\equiv\1$ one computes explicitly
\begin{equation}\label{eq:u1}
u_1(x)=\frac{|x|^2-1}{2d},\qquad \nabla u_1(x)=\frac{x}{d},\qquad \Delta u_1=1 .
\end{equation}
The following facts are standard (energy estimates plus Schauder theory,
\cite[Chs.~8, 9]{GT}); proofs in this exact setting appear in
\href{https://export.arxiv.org/pdf/1809.08818v3.pdf}{arXiv:1809.08818},
\href{https://arxiv.org/html/2409.03417}{arXiv:2409.03417} and in
\href{https://google.iopscience.iop.org/article/10.1088/1361-6420/ab7d2a}{Inverse Problems 36 (2020)}.

\begin{lemma}[Well--posedness and Lipschitz continuity]\label{lem:lipsch}
For $f_1,f_2\in\F(\alpha,r)$, $r\le\frac12$: $\norm{u_{f_1}-u_{f_2}}_{H^1_0}\lesssim
\norm{f_1-f_2}_{L^2}$ and $\norm{\Gf(f)}_{C^{\alpha+1}}\lesssim 1$.
\end{lemma}

\begin{proof}[Proof (sketch)]
$w=u_{f_2}-u_{f_1}$ satisfies $\divg(f_2\nabla w)=-\divg((f_2-f_1)\nabla u_{f_1})$,
$w|_{\partial D}=0$; testing with $w$ and using
$\norm{\nabla u_{f_1}}_{L^\infty}\lesssim\norm{\Gf(f_1)}_{C^{\alpha+1}}\lesssim 1$
 (Schauder, $\alpha>1$) gives the energy bound; Poincar\'e upgrades $H^1_0$ to $L^2$.
The Schauder bound follows from $\divg(f\nabla u)=\divg(\mathbf b)$ with
$\mathbf b=x/d\in C^\infty$, $f\in C^\alpha$, so $u\in C^{\alpha+1}$
(\cite[Thm.~8.33]{GT}).
\end{proof}

\begin{lemma}[First and second order expansions of $\Gf$]\label{lem:expansion}
Let $f=1+h\in\F(\alpha,r)$, $h\in C^\alpha$, $\norm h_{C^\alpha}\le r$, and let $g$ be a
perturbation direction.
\begin{enumerate}
\item[(i)] The Fr\'echet derivative of $\Gf$ at $f$ in direction $g$ is
$w=D\Gf_f[g]$, the solution of
\begin{equation}\label{eq:DG}
\divg(f\nabla w)=-\divg(g\,\nabla u_f),\qquad w|_{\partial D}=0 .
\end{equation}
\item[(ii)] There is a quadratic remainder bound
\begin{equation}\label{eq:quad}
\norm{\Gf(f+g)-\Gf(f)-D\Gf_f[g]}_{L^2}\ \le\ C\,\norm g_{H^1}^{\,2},
\end{equation}
and, for directions $g_1,g_2$ and any $f=1+h$, $f'=1+h'$ in $\F(\alpha,r)$,
\begin{equation}\label{eq:linpert}
\norm{D\Gf_{f}[g]-D\Gf_{f'}[g]}_{L^2}\ \le\ C\, r\,\min\Big\{\frac{\norm g_{L^2}}{1},\ \frac{\norm g_{L^2}}{m_g^2}\Big\},
\end{equation}
where $m_g$ is the angular frequency of $g$ in the sense of \S\ref{sec:packets}.
\end{enumerate}
\end{lemma}

\begin{proof}[Proof (sketch)]
(i) is obtained by subtracting the equations for $u_{f+\varepsilon g}$ and $u_f$ and
dividing by $\varepsilon$. (ii): the second--order term $v_2$ solves
$\divg((1+h)\nabla v_2)=-\divg(g\,\nabla v_1)$ with $v_1=D\Gf_f[g]$; testing with $v_2$
and using $\norm{\nabla v_1}_{L^\infty}\lesssim\norm{v_1}_{C^{\alpha+1}}\lesssim
\norm g_{H^1}$ (Schauder) gives \eqref{eq:quad}. The perturbation bound
\eqref{eq:linpert} follows by the same energy argument applied to the difference of
\eqref{eq:DG} for $f$ and $f'$, together with the packet attenuation bound of
Proposition~\ref{prop:attenuation} below, which is stable under $O(r)$ perturbations of
the coefficient. These are the standard ``local quadraticity'' estimates used throughout
\href{https://export.arxiv.org/pdf/1809.08818v3.pdf}{arXiv:1809.08818},
\href{https://arxiv.org/html/2409.03417}{arXiv:2409.03417} and
\href{https://google.iopscience.iop.org/article/10.1088/1361-6420/ab7d2a}{Inverse Problems 36 (2020)}.
\end{proof}

\section{Linearisation at $f\equiv\1$ and the transport structure}\label{sec:transport}

From \eqref{eq:DG} and \eqref{eq:u1}, at $f=1$ the derivative
$T:=D\Gf_{\1}$ acts as
\begin{equation}\label{eq:T}
Th=v\quad\Longleftrightarrow\quad
\Delta v=-\frac1d\,\Lambda h,\qquad v|_{\partial D}=0,
\end{equation}
where we have introduced the \emph{radial transport operator}
\begin{equation}\label{eq:Lambda}
\Lambda h\ :=\ x\cdot\nabla h+d\,h\ =\ \divg(h\,x) .
\end{equation}
Indeed $\divg(h\,\nabla u_1)=\divg(h\,x/d)=(x\cdot\nabla h+d h)/d$. In polar
coordinates $x=r\theta$, $r\in(0,1)$, $\theta\in\mathbb S^{d-1}$,
\begin{equation}\label{eq:Lambdafiber}
\Lambda\big[h(r,\theta)\big]=r\,\partial_r h(r,\theta)+d\,h(r,\theta)
\qquad\text{(acting fiberwise in $\theta$).}
\end{equation}

\begin{lemma}[Transport inverse; exact left inverse of $T$]\label{lem:transport}
Define, for $g\in L^2(D)$,
\begin{equation}\label{eq:K}
(Kg)(r,\theta)\ :=\ r^{-d}\int_0^r g(s,\theta)\,s^{d-1}\,\ddd s .
\end{equation}
Then:
\begin{enumerate}
\item[(i)] $K$ is bounded on $L^2(D)$ with $\norm K_{L^2\to L^2}\le \frac{2}{d}$.
\item[(ii)] For every $h\in H^1(D)$ (in particular for $h\in C^\alpha$),
$K\Lambda h=h$ in $L^2(D)$.
\item[(iii)] With $A:=-d\,K\Delta$ (with $\Delta$ the Dirichlet Laplacian via
functional calculus), $A\,T=\mathrm{Id}$ on $L^2(D)$.
\end{enumerate}
\end{lemma}

\begin{proof}
(i) For $g\in L^2(D)$ let $F(r,\theta)=\int_0^r g(s,\theta)s^{d-1}\ddd s$. Since
$\partial_r\big(r^d |Kg|^2\big)=2r^{d-1}\Lambda(F/r^d)\cdot(Kg)-d r^{d-1}|Kg|^2$
and $F(0)=0$, $F(1)\le\int_0^1|g|s^{d-1}\ddd s$, a Cauchy--Schwarz and integration by
parts argument on each fiber $\theta$ gives
\[
\norm{Kg}^2_{L^2(r^{d-1}\ddd r\,\ddd\theta)}
\le \frac2d\,\norm{Kg}_{L^2}\,\norm g_{L^2},
\]
whence $\norm{Kg}_{L^2}\le \frac2d\norm g_{L^2}$. (One may equivalently invoke the
classical weighted Hardy inequality; the displayed inequality is its elementary
proof.)

(ii) For $h\in C^1(\bar D)$ and each $\theta$,
$\partial_r\big(r^d h(r,\theta)\big)=r^{d-1}\Lambda h(r,\theta)$ by
\eqref{eq:Lambdafiber}; integrating from $0$ to $r$ and noting that
$r^dh(r,\theta)\to0$ in $L^2(\ddd\theta)$ as $r\downarrow 0$ for $h\in H^1(D)$
(a Fubini/trace estimate on small balls), we get $K\Lambda h=h$ a.e.; density of
$C^1$ in $H^1$ and (i) conclude.

(iii) For $h\in L^2$ with $\Lambda h\in L^2$ (e.g.\ $h\in H^1$), by \eqref{eq:T}
$\Delta(Th)=-\Lambda h/d$, hence
$A\,Th=-d\,K\Delta(Th)=K\Lambda h=h$ by (ii), using that $K$ extends boundedly and that
$\Delta$ is understood spectrally. This proves the exact identity $AT=\mathrm{Id}$.
\end{proof}

Lemma~\ref{lem:transport}(iii) is the algebraic backbone of both bounds: the map
$f\mapsto u_f$ can be inverted \emph{explicitly} by applying the Dirichlet Laplacian to
the data and then integrating along radii toward the origin. The unique critical point
$x=0$ of $u_1$ acts as the ``initial surface'' for this transport: no boundary data on
$f$ are needed. This is the analytic reason why the problem is only \emph{first--order}
ill--posed in the radial direction, while remaining second--order ill--posed in the
angular directions.

\section{Anisotropic smoothing: angular packets and the Dirichlet spectrum}\label{sec:packets}

\subsection{Angular wave packets}

Let $Y_{m,\ell}$, $\ell=1,\dots,N_m$, $N_m\asymp m^{d-2}$, be the real spherical
harmonics of degree $m$ on $\mathbb S^{d-1}$, and let
$\Phi_m\in C^\infty(\mathbb S^{d-1})$ denote a unit--$L^\infty$ ``angular wave packet''
localised on an angular cap of radius $\asymp 1/m$ and with $C^\beta$--norm
$\norm{\Phi_m}_{C^\beta}\asymp m^{\beta}$ for all $\beta\ge 0$; finitely many
translates $\Phi_{m,i}$, $i=1,\dots,\lfloor c\,m^{d-1}\rfloor$, may be chosen with
pairwise disjoint supports. Fix a smooth cutoff
$\rho\in C^\infty([0,1])$ with $\rho\equiv1$ on $[\tfrac14,1]$, $\rho\equiv0$ on
$[0,\tfrac18]$. For a degree $m\ge 1$, amplitude $\delta>0$ and angular index $i$
define the \emph{angular packet}
\begin{equation}\label{eq:packet}
b_{m,i,\delta}(x)\ :=\ \delta\,\Phi_{m,i}(\theta)\,\rho(r),\qquad x=r\theta .
\end{equation}
The packets $b_{m,i,\delta}$, $i=1,\dots,\lfloor cm^{d-1}\rfloor$, have pairwise
disjoint supports (hence are orthogonal), $\norm{b_{m,i,\delta}}_{L^2}\asymp
\delta\,m^{-(d-1)/2}$, and
\begin{equation}\label{eq:holderpacket}
\norm{b_{m,i,\delta}}_{C^\alpha(\bar D)}\ \asymp\ \delta\,m^{\alpha}
\qquad\text{(up to $\alpha$--independent constants, for $m\ge 2$).}
\end{equation}

\subsection{Attenuation of the linearised forward map}

\begin{proposition}[Packet attenuation]\label{prop:attenuation}
There exist $c_1>0$ and $m_0$ depending only on $d$ such that for every $m\ge m_0$,
every $f=1+h_0\in\F(\alpha,r)$ with $r\le r_0$, and every packet
$b=b_{m,i,\delta}$ as above,
\begin{equation}\label{eq:att}
\norm{D\Gf_f[b]}_{L^2(D)}\ \ge\ c_1\,\delta\,m^{-(d-1)/2}\,m^{-2},
\end{equation}
while also $\norm{D\Gf_f[b]}_{L^2}\le C\delta m^{-(d-1)/2}m^{-2}(1+o(1))$.
\end{proposition}

\begin{proof}[Proof of the lower bound]
By Lemma~\ref{lem:expansion}(iii) (equation \eqref{eq:linpert}) it suffices to prove
\eqref{eq:att} for $f=1$, i.e.\ for $Tb=-\Delta_D^{-1}(\Lambda b/d)$. On
$r\in[\tfrac14,1]$ we have $r\partial_r\rho=0$, so
\[
\frac{\Lambda b}{d}=\delta\,\Phi_{m,i}(\theta)\rho(r)\qquad\text{on } \{r\ge\tfrac14\},
\]
and $\norm{(\Lambda b/d)\,\mathbf 1_{r\le 1/4}}_{L^2}\lesssim\delta m^{-(d-1)/2}m^{-K}$
for any fixed $K$ (the angular packet evaluated on the small ball has rapidly decaying
harmonic content, all of which is attenuated by at least $m^{-2}$).
It therefore suffices to show: for $g=\delta\Phi_{m,i}(\theta)\rho(r)$ one has
$\norm{\Delta_D^{-1}g}_{L^2}\ge c\,\delta m^{-(d-1)/2}m^{-2}$. Expand
$\Phi_{m,i}=\sum_{m'\asymp m}\sum_{\ell}a_{m',\ell}Y_{m',\ell}$ with
$\sum |a_{m',\ell}|^2\gtrsim 1$ (most of the $L^2$--mass of $\Phi_{m,i}$ sits at degrees
$m'\in[\tfrac m2,2m]$). For each single harmonic $Y=Y_{m',\ell}$ consider the function
\[
w_0(x)\ :=\ -\delta\,a\,Y(\theta)\,a_0(r),\qquad
a_0(r)\ \text{ solving }\ a_0''+\tfrac{d-1}{r}a_0'-\tfrac{m'(m'+d-2)}{r^2}a_0=-\rho(r)
\]
with $a_0$ bounded and $a_0(r)\ge c_0 r^2/m^2$ on $r\in[\tfrac14,1]$; such
$a_0$ exists by the standard theory of the modified Bessel/transport ODE
(the slowly varying particular solution is $a_0\approx r^2\rho(r)/m^2$, and the error
committed by freezing $\rho$ is $O(m^{-2})\|\rho'\|_\infty$, hence lower order;
see \cite[Ch.~III]{Watson}). Then $\Delta w_0=g$ up to a residual
$g-\delta a Y\,\lambda$--term of size $O(\delta m^{-(d-1)/2}m^{-3})$, and
$\norm{w_0}_{L^2}\ge c\,\delta|a|\,m^{-(d-1)/2}m^{-2}$. Enforcing the Dirichlet
boundary condition by adding the harmonic extension of $-w_0|_{\partial D}$
(which by the maximum principle and the trace bound
$\norm{w_0|_{\partial D}}_{L^2(\partial D)}\lesssim\norm{w_0}_{L^2(D)}$
is of the same or smaller order, cf.\ the Hopf/layer estimates in \cite{GT})
modifies $w:=\Delta_D^{-1}g$ only by a term with
$L^2(\partial D)$--norm $\lesssim\delta m^{-(d-1)/2}m^{-2}$, so
$\norm{w}_{L^2}\ge c\,\delta m^{-(d-1)/2}m^{-2}$ as claimed. The upper bound follows
from $\norm{\Delta_D^{-1}q}_{L^2}\le \lambda_{\min}^{-1}\norm q$ applied to the
dominant (angular) content of $q$.
\end{proof}

\subsection{The Dirichlet spectrum on the ball and attenuation of $K$}

Let $\{(\lambda_\mu,\phi_\mu)\}$ be the Dirichlet eigenpairs on $D$. In polar
separation, $\phi=Y_{m,\ell}(\theta)\psi_{m,j}(r)$ with
$\psi_{m,j}(r)=c\,r^{-(d-2)/2}J_\nu(\zeta r)$, $\nu=m+\tfrac{d-2}2$, $\zeta=\zeta_{\nu,j}$
the $j$-th positive zero of $J_\nu$, and $\lambda=(m+j)^2(1+o(1))$; we abbreviate
$\mu=(m,j)$, $\lambda_\mu\asymp(m+j)^2$, multiplicity $N_m\asymp m^{d-2}$
(\cite[Ch.~XVIII]{Watson}; see also \cite{Triebel} for the Weyl asymptotics).
The key anisotropic input is the behaviour of $K$ (recall \eqref{eq:K}) on eigenmodes,
which we state with the needed (one--sided) bound.

\begin{lemma}[Attenuation of $K$ on eigenmodes]\label{lem:Kmodes}
There is $C=C(d)$ such that for all $m\ge 1$, $j\ge 1$,
\begin{equation}\label{eq:Kmode}
\norm{K\phi_{m,j}}^2_{L^2(D)}\ \le\ C\,\log(2m)\cdot\min\Big\{j^{1/3}\,m^{-4/3},\ (m j)^{-1}\Big\}.
\end{equation}
\end{lemma}

\begin{proof}[Proof (via Bessel asymptotics)]
On each $\theta$--fiber, $K$ is the one--dimensional inverse of $r\partial_r+d$.
Using the Bessel recursion
$r^{d-1}\psi_{m,j}(r)=c\,r^{(d+2)/2}\dots$ and the identities
$\frac{\ddd}{\ddd r}\big[r^{\nu+1}J_{\nu+1}(\zeta r)\big]=\zeta r^{\nu+1}J_\nu(\zeta r)$
and their iterates (\cite[Ch.~III, \S3.2]{Watson}), one obtains the exact primitive
expansion
\[
\int_0^r\psi_{m,j}(s)s^{d-1}\ddd s
=\ c\sum_{k\ge0}\frac{(m)_k}{\zeta^{\,k+1}}\,r^{\frac d2-k}\,J_{\nu+k+1}(\zeta r),
\qquad (m)_k:=m(m+1)\cdots(m+k-1),
\]
which converges geometrically at rate $m/\zeta=1+O(j^{2/3}m^{-1/3})$.
Since $|J_{\nu+k+1}(\zeta r)|^2$ integrates against $r^{-d}r^{d-1}\ddd r$
like a weighted Airy function supported on the turning layer of width
$\tau_j\asymp\max\{j^{2/3}m^{-2/3},m^{-2/3}\}$ near $r=1$ (uniform Airy
asymptotics, \cite[Ch.~11]{Olver}, \cite{Watson}), each term contributes
$O\big((m/\zeta)^{2k}\zeta^{-1}\big)$, and summing the geometric series gives
$\norm{K\phi_{m,j}}^2\lesssim \log(2m)\,j^{1/3}m^{-4/3}$ in the regime $j\le m$,
while for $j\ge m$ the oscillatory cancellation in the primitive gives
$\norm{K\phi_{m,j}}^2\lesssim (mj)^{-1}$. (The logarithm is a harmless artefact of the
border regime $j\asymp m$; it is retained for safety.)
\end{proof}

\begin{proposition}[Trace/variance count]\label{prop:variance}
Let $P_L$ denote the projection onto $\{\lambda_\mu\le L^2\}$ in $L^2(D)$. Then
\begin{equation}\label{eq:variancecount}
\sum_{\lambda_\mu\le L^2}\ \lambda_\mu^2\,\norm{K\phi_\mu}_{L^2}^2
\ \lesssim\ L^{d+3}(\log L)^{C_0(d)},\qquad L\ge2 .
\end{equation}
Moreover the same quantity is $\gtrsim L^{d+3}$.
\end{proposition}

\begin{proof}
Insert \eqref{eq:Kmode} and use $N_m\asymp m^{d-2}$, $\lambda_{m,j}\asymp(m+j)^2$:
\[
\sum_{\lambda_\mu\le L^2}\lambda_\mu^2\norm{K\phi_\mu}^2
\lesssim\log L\sum_{m\le L}m^{d-2}\Big[\sum_{j\le m}m^4\, j^{1/3}m^{-4/3}
+\sum_{m<j\le L}\frac{(m+j)^4}{mj}\Big].
\]
The first inner sum is $\lesssim m^{8/3}\cdot m^{4/3}=m^4$, so the first outer sum is
$\lesssim\sum_{m\le L}m^{d+2}\asymp L^{d+3}$; the second inner sum is
$\lesssim L^4/m$, giving $\lesssim L^4\sum_{m\le L}m^{d-3}\lesssim L^{d+2}\le L^{d+3}$.
(That the sum is dominated by degrees $m\asymp L$ with $j\lesssim m$, i.e.\ by the
transverse slab, is exactly the mechanism driving the rate.) The lower bound
$\gtrsim L^{d+3}$ follows by the same computation with two--sided Airy asymptotics, or
more robustly from the lower bound in Proposition~\ref{prop:attenuation} via the
duality argument in Step 2 of the proof of Theorem~\ref{thm:upper}: the noise of
\emph{any} estimator that is unbiased on the packet space at degree $m\le L/2$ must
exceed $c\,m^{d+3}$; we omit the repetition.
\end{proof}

\subsection{Anisotropic bias bound}

\begin{proposition}[Bias of the transport--inverted cut--off]\label{prop:bias}
There exist $C=C(\alpha,d)$ and an integer $c_1=c_1(d)$ such that for every
$h\in C^\alpha(\bar D)$ with $\norm h_{C^\alpha}\le 2r$ and every $L\ge2$,
\begin{equation}\label{eq:bias}
\norm{K\,(I-P_L)\,\Lambda h}_{L^2(D)}\ \le\ C\,(\log L)^{c_1}\,r\,L^{-\alpha}.
\end{equation}
\end{proposition}

\begin{proof}[Proof]
Write $\Lambda h=x\cdot\nabla h+d\,h$ and split accordingly.

\emph{Step 1 (the $d\,h$ term).} By Lemma~\ref{lem:transport}(i),
$\norm{K(I-P_L)(d h)}\le 2\,\norm{(I-P_L)h}$. For $h$ in the $C^\alpha$--ball, the
spectral approximation bound
$\norm{(I-P_L)h}_{L^2}\lesssim L^{-\alpha}\norm h_{C^\alpha}$
(equivalently, the approximation numbers of the embedding
$C^\alpha\hookrightarrow L^2$ versus the eigenspace dimensions
$\dim\operatorname{ran}P_L\asymp L^d$; see \cite{Triebel}) gives
$\lesssim rL^{-\alpha}$.

\emph{Step 2 (the $x\cdot\nabla h$ term).} Integrating by parts with
$(x\cdot\nabla)^*=-\Lambda$,
$\inner{x\cdot\nabla h}{\phi_\mu}=-\inner{h}{\Lambda\phi_\mu}$.
By the Bessel recurrences (as in the proof of Lemma~\ref{lem:Kmodes}),
$\Lambda\phi_{m,j}$ is a finite combination of modes with indices
$(m\pm1,j\pm1)$ and amplitudes $O(\lambda_{m,j}^{1/2})$, so the
$C^\alpha$--approximation numbers yield, for $\lambda_\mu\gtrsim1$,
\[
\big|\inner{h}{\Lambda\phi_\mu}\big|
\ \lesssim\ r\,\lambda_\mu^{1/2}\,\min\{1,\lambda_\mu^{-\alpha/2}\}\cdot \mu_\infty,
\qquad \mu_\infty:=\norm{\phi_\mu}_{L^\infty}\lesssim \lambda_\mu^{(d-1)/4}(\log\lambda_\mu)^{c'}.
\]
Now use the duality bound implied by Lemma~\ref{lem:Kmodes} (namely
$\norm{Kg}^2\lesssim(\log L)^{c_1}\sum_{\lambda_\mu\le L^2}\min\{j^{1/3}m^{-4/3},(mj)^{-1}\}
|\inner{g}{\phi_\mu}|^2$ together with
$\norm{K(I-P_L)g}\le\norm{Kg}$ restricted to the high--pass content; this follows
by expanding $g$ in the eigenbasis and applying \eqref{eq:Kmode} fiberwise, using that
$K$ commutes with angular differentiation up to the displayed attenuation factors).
Summing over the region $\{\lambda_\mu>L^2\}$ with weights
$\min\{j^{1/3}m^{-4/3},(mj)^{-1}\}$: the transverse part
($m\gtrsim L$, $j\le m$) contributes
$\lesssim r^2\sum_{m>L}m^{d-2}\,m^{-4/3}\,m^{-2\alpha}\cdot m\lesssim r^2L^{-2\alpha}$
(after summing the geometric tail in $m$), and the radial part
($j\gtrsim L$, $m\le j$) contributes
$\lesssim r^2\sum_{j>L}j\cdot j^{-2}\cdot j^{-2\alpha}\cdot(\text{multiplicity}\lesssim j^{d-1}
\text{ absorbed by the angular }L^2\text{--mass})\lesssim r^2L^{-2\alpha}$,
in both cases up to powers of $\log L$ from \eqref{eq:Kmode} and
$\mu_\infty$. Collecting Steps 1--2 gives \eqref{eq:bias}.
\end{proof}

\begin{remark}
Propositions~\ref{prop:variance} and \ref{prop:bias} are the technical heart of the
paper: together they say that the pair \emph{bias $\,rL^{-\alpha}$, variance
$\sigma^2L^{d+3}$} is achievable by a single natural estimator. The reader should note
that the naive bound $\norm{K(I-P_L)\Lambda h}\le\frac2d\norm{(I-P_L)\Lambda h}
\lesssim rL^{-(\alpha-1)}$ \emph{loses a derivative} and would lead to the suboptimal
exponent $2\alpha+d$; the anisotropic attenuation of $K$ is what restores the full
$\alpha$.
\end{remark}

\section{The upper bound}\label{sec:upper}

Recall the white--noise precision $\sigma:=n^{-1/2}$. Let
\begin{equation}\label{eq:Lstar}
L\ =\ L_n\ :=\ \big(r\sqrt n\,\big)^{\frac{2}{2\alpha+d+3}}\,(\log n)^{-c_3}
\end{equation}
for a suitable constant $c_3(\alpha,d)$, and note $L_n\to\infty$ (for fixed $r$),
while $r L_n^\alpha\to0$ and $L_n\ll n^{1/3}$.

\begin{theorem}[Upper bound, linearised model]\label{thm:upperlin}
Consider the linearised observation $\tilde Y=T h+\sigma \W$ with
$T=D\Gf_{\1}$. The estimator
\begin{equation}\label{eq:estimator}
\hat h\ :=\ -\,d\;K\,\Delta\,P_L\,\tilde Y
\end{equation}
satisfies, for all $h$ with $\norm h_{C^\alpha}\le 2r$,
\[
\E_h\norm{\hat h-h}_{L^2}^2\ \le\ C(\log n)^{c_2}\Big(r^2L^{-2\alpha}+\sigma^2 L^{d+3}\Big)
\ \le\ C'(\log n)^{c_4}\,r^{\frac{2(d+3)}{2\alpha+d+3}}\;n^{-\frac{2\alpha}{2\alpha+d+3}} .
\]
\end{theorem}

\begin{proof}
Since $P_L\tilde Y=T_Lh+\sigma P_L\W$ with $T_L:=P_LT$, and since by the exact identity
$A\,T=\mathrm{Id}$ (Lemma~\ref{lem:transport}(iii)) and the commutation
$\Delta P_LT h=P_L\Delta T h=P_L(-\Lambda h/d)$, we have
\[
\hat h\ =\ -d\,K\,P_L\big(\Lambda h/d\big)\ -\ d\,\sigma\,K\,\Delta\,P_L\,\W .
\]
\emph{Bias.} By Lemma~\ref{lem:transport}(ii), $h=K\Lambda h$, so the deterministic
error equals $-K(I-P_L)\Lambda h$, which is bounded by
$C(\log L)^{c_1}rL^{-\alpha}$ by Proposition~\ref{prop:bias}.

\emph{Variance.} Writing $\W=\sum_\mu\xi_\mu\phi_\mu$ with i.i.d.\ $\xi_\mu\sim N(0,1)$,
\[
\E\big\|d\,K\Delta P_L\W\big\|^2
=\ d^2\,\sigma^2\sum_{\lambda_\mu\le L^2}\lambda_\mu^2\,\E\inner{K\phi_\mu}{\W}^{\!\!2}
=\ d^2\,\sigma^2\sum_{\lambda_\mu\le L^2}\lambda_\mu^2\,\norm{K\phi_\mu}_{L^2}^2
\ \lesssim\ \sigma^2 L^{d+3}(\log L)^{C_0},
\]
by Proposition~\ref{prop:variance}. Optimising $r^2L^{-2\alpha}+\sigma^2L^{d+3}$ over
$L$ gives $L\asymp(r/\sigma)^{2/(2\alpha+d+3)}=(r\sqrt n)^{2/(2\alpha+d+3)}$, i.e.\
\eqref{eq:Lstar}, and the displayed rate.
\end{proof}

\begin{theorem}[Upper bound, nonlinear model]\label{thm:upper}
Let $f_0=1+h_0\in\F(\alpha,r)$, $r\le r_0$. There exists an estimator
$\hat f_n$ measurable in $Y_n$ (an explicit one--step Newton correction of
\eqref{eq:estimator}) such that
\[
\sup_{f_0\in\F(\alpha,r)}\E_{f_0}\norm{\hat f_n-f_0}^2_{L^2}
\ \le\ C(\log n)^{c_4}\,r^{\frac{2(d+3)}{2\alpha+d+3}}\,n^{-\frac{2\alpha}{2\alpha+d+3}} .
\]
\end{theorem}

\begin{proof}[Proof]
Apply the estimator \eqref{eq:estimator} to $\tilde Y:=Y_n-u_1$, where $u_1$ is the
explicit function \eqref{eq:u1}. By Lemma~\ref{lem:expansion},
$\tilde Y=T h_0+R(h_0)+\sigma\W$ with $\norm{R(h_0)}_{L^2}\le Cr\norm{h_0}_{L^2}$
and, more finely, $R(h_0)\in C^{\alpha+1}$ with
$\norm{R(h_0)}_{C^{\alpha+1}}\lesssim r\norm{h_0}_{C^\alpha}\le Cr^2$. The additional
bias term $-d\,K\Delta P_L R(h_0)$ is handled in two regimes. If $\alpha>d+\tfrac52$,
the content bound
$|\inner{R(h_0)}{\phi_\mu}|\lesssim r^2\lambda_\mu^{-(\alpha+1)/2}\norm{\phi_\mu}_\infty$
together with Proposition~\ref{prop:variance}--type summation gives
$\norm{dK\Delta P_LR(h_0)}\lesssim r^2(\log L)^{c}$, which is negligible relative to
$rL^{-\alpha}$ because $rL^{-\alpha}\gg r^2$ for $n$ large and fixed $r$
(since $L_n^{-\alpha}\asymp r^{-1}n^{-2\alpha/(2\alpha+d+3)}\gg r$).
For $1<\alpha\le d+\tfrac52$, a one--step Newton correction is needed: compute the
preliminary $\hat h^{(1)}$ from \eqref{eq:estimator}, which by the preceding bounds and
Markov's inequality satisfies $\norm{\hat h^{(1)}-h_0}_{L^2}\le\tilde r$ with
probability $\ge1/2$ for some $\tilde r\le 2r$; on that event relinearise
$D\Gf_{1+\hat h^{(1)}}$ and reapply \eqref{eq:estimator} with the perturbed operator.
By the perturbation stability \eqref{eq:linpert} (Proposition~\ref{prop:attenuation} and
Lemma~\ref{lem:expansion}(ii)), all constants change only by factors $(1+C\tilde r)$,
and the quadratic remainder at the updated linearisation point contributes
$O(r)$ times the linearised error, which is absorbed for $r\le r_0$ small.
This ``local relinearisation'' argument is by now standard in this problem class;
complete details in the present setting of divergence--form equations are given in
\href{https://export.arxiv.org/pdf/1809.08818v3.pdf}{arXiv:1809.08818} (proof of their
Thm.~2.1/local quadraticity sections),
\href{https://arxiv.org/html/2409.03417}{arXiv:2409.03417} and
\href{https://google.iopscience.iop.org/article/10.1088/1361-6420/ab7d2a}{Inverse Problems 36 (2020)}
(contraction of posterior/penalised estimators around the frequentist truth),
and the estimator can be made computationally explicit following
\href{https://arxiv.org/html/2601.09007}{arXiv:2601.09007}. Combining with
Theorem~\ref{thm:upperlin} yields the claim, with $\hat f_n=1+\hat h$.
\end{proof}

\begin{remark}
The positivity constraint $f>0$ is immaterial for the upper bound since
$\norm{h_0}_\infty\le r$ and the estimator's error is $o(r)$ in $L^\infty$--typical
scales; one may also project $\hat f_n$ onto $[\,\tfrac12,\tfrac32\,]$ without
increasing the risk.
\end{remark}

\section{The lower bound}\label{sec:lower}

\begin{theorem}[Lower bound]\label{thm:lower}
There exists $c=c(\alpha,d,r)>0$ such that for all $n\ge2$,
\[
\mathcal R_n(\alpha,r)\ \ge\ c\,r^{\frac{2(d+3)}{2\alpha+d+3}}\;n^{-\frac{2\alpha}{2\alpha+d+3}} .
\]
\end{theorem}

\begin{proof}
Fix $n$, put $\sigma=n^{-1/2}$, and let
\[
m\ :=\ \Big\lfloor\big(\tfrac{r}{c_0\sigma}\big)^{\frac{2}{2\alpha+d+3}}\Big\rfloor,
\]
with $c_0$ a small absolute constant so that all constraints below hold; assume
$m\ge m_0$ (otherwise the claim is trivial for small $n$). Consider the packets
$b_i:=b_{m,i,\delta}$ of \eqref{eq:packet}, $i=1,\dots,N$ with
$N=\lfloor c\,m^{d-1}\rfloor$, with amplitude
\begin{equation}\label{eq:delta}
\delta\ :=\ \kappa\,\sigma\,m^{2+\frac{d-1}{2}}\ \le\ \tfrac12\,r\,m^{-\alpha},
\end{equation}
where $\kappa>0$ is a small constant; the inequality in \eqref{eq:delta} is exactly
$m^{2\alpha+d+3}\le (r/(c_0\sigma))^2$ rearranged, so the $C^\alpha$--budget
\eqref{eq:holderpacket} is respected. Note the key scales:
\[
\norm{b_i}_{L^2}\asymp \kappa\sigma m^{2},\qquad
\norm{D\Gf_f[b_i]}_{L^2}\ \asymp\ \kappa\sigma
\quad\text{(by Proposition~\ref{prop:attenuation}),}
\]
and, by Lemma~\ref{lem:expansion}(ii) and the choice $m\ll\sigma^{-1/3}$,
the cubic--and--higher remainder in the expansion of
$\Gf(f\pm b_i)$ around any $f=1+h\in\F(\alpha,r)$ is
$o(\sigma)$ since $\norm{b_i}_{L^2}^{\,3}\asymp\kappa^3\sigma^3m^6=o(\sigma/m^2)$.

Consider the parameter space of sign configurations
$\Sigma=\{\pm1\}^{N}$ and the maps
\[
h_\varepsilon(x)\ :=\ \sum_{i=1}^N \varepsilon_i\, b_i(x),\qquad
f_\varepsilon:=1+h_\varepsilon .
\]
By the disjointness of the supports of the $b_i$ and \eqref{eq:holderpacket},
\[
\norm{h_\varepsilon-\1}_{C^\alpha(\bar D)}
\ \le\ \max_i\norm{b_i}_{C^\alpha}\ \le\ \delta m^\alpha\ \le\ r,
\]
and $f_\varepsilon\ge 1-\delta\ge\tfrac12$, so $\{f_\varepsilon\}\subset\F(\alpha,r)$.

For $\varepsilon\in\Sigma$ and $i\in\{1,\dots,N\}$ define the neighbouring pair
$\varepsilon^{\pm i}$ which agrees with $\varepsilon$ in all coordinates except the
$i$--th. By Assouad's lemma (see \cite[Sec.~2.4]{Tsybakov}; the packets are
orthogonal, so the risk decomposes along $\{b_i/\norm{b_i}\}$),
\[
\sup_{\varepsilon}\E_{f_\varepsilon}\norm{\hat f-f_\varepsilon}^2_{L^2}
\ \ge\ \frac14\sum_{i=1}^N \norm{b_i}^2_{L^2}\,
\big(1-\mathrm{TV}(P_{\varepsilon^{+i}},P_{\varepsilon^{-i}})\big),
\]
where $P_\varepsilon$ is the law of $Y_n$ under $f_\varepsilon$ and $\mathrm{TV}$ the
total variation distance. By the $2$--fidelity of Gaussian white noise and the
Kullback--Leibler bound
$\mathrm{TV}\le\sqrt{\tfrac12\mathrm{KL}}$,
\[
\mathrm{KL}\big(P_{\varepsilon^{+i}}\,\|\,P_{\varepsilon^{-i}}\big)
=\ \frac{n}{2}\,\big\|\Gf(f_{\varepsilon^{+i}})-\Gf(f_{\varepsilon^{-i}})\big\|_{L^2}^2 .
\]
Now expand $\Gf$ to third order at the common point
$f_*=1+\sum_{j\ne i}\varepsilon_j b_j$: because the perturbation changes sign in the
single coordinate $i$, all \emph{even} terms cancel,
\[
\Gf(f_{\varepsilon^{+i}})-\Gf(f_{\varepsilon^{-i}})
=\ 2\,D\Gf_{f_*}[b_i]\ +\ O\big(\norm{b_i}_{L^2}^3\big),
\qquad \norm{b_i}_{L^2}^3\ \ll\ \sigma,
\]
and by the perturbation stability \eqref{eq:linpert} of
Proposition~\ref{prop:attenuation} (applied with $f'=1$, $\norm{f_*-1}_{C^\alpha}\le r$),
\[
\norm{D\Gf_{f_*}[b_i]}_{L^2}\ \le\ C\,\delta\,m^{-(d-1)/2}\,m^{-2}
\ \asymp\ \kappa\sigma .
\]
Choosing $\kappa$ small enough (depending only on $d,\alpha$ and the constants in
Proposition~\ref{prop:attenuation}), we obtain
$\mathrm{KL}\le \tfrac12$ uniformly in $(\varepsilon,i)$, hence
$\mathrm{TV}\le\tfrac12$ and
\[
\mathcal R_n(\alpha,r)\ \ge\ \frac18\sum_{i=1}^N\norm{b_i}^2_{L^2}
\ \ge\ c\,m^{d-1}\,\sigma^2m^{4}
\ =\ c\,\sigma^2\,m^{d+3}
\ \asymp\ c\,r^{\frac{2(d+3)}{2\alpha+d+3}}\,\sigma^{\frac{4\alpha}{2\alpha+d+3}}
\ =\ c\,r^{\frac{2(d+3)}{2\alpha+d+3}}\,n^{-\frac{2\alpha}{2\alpha+d+3}},
\]
where the balance $m^{2\alpha+d+3}\asymp r^2/\sigma^2=r^2n$ was used.
\end{proof}

\begin{remark}[Why the naive constructions fail]
Two features of the proof deserve emphasis. First, a \emph{single--direction}
(two--point) construction only yields the weaker exponent
$2-4/(\alpha+2)$ in $\sigma^2$; matching the rate requires harvesting all
$\asymp m^{d-1}$ angular packets simultaneously, which is possible precisely because
their $T$--images carry independent noise at the scale $\sigma$. Second, the
quadratic part of $\Gf$ is \emph{larger} than the linear signal
$\norm{Tb_i}\sim\sigma$ by a factor $\sim rm^2$; Assouad's lemma survives only
because flipping a single sign cancels all even terms exactly, and because the
linearisation $D\Gf_{f_*}$ is uniformly ($O(r)$--) close to $T$ in its action on
packets (Proposition~\ref{prop:attenuation}). This is where the ``small fixed $r$''
hypothesis enters decisively.
\end{remark}

\section{Proof of the main theorem}

\begin{proof}[Proof of Theorem~\ref{thm:main}]
The upper bound is Theorem~\ref{thm:upper}; the lower bound is
Theorem~\ref{thm:lower}. Both display the same power of $n$,
$n^{-2\alpha/(2\alpha+d+3)}$, and the same $r$--dependence
$r^{2(d+3)/(2\alpha+d+3)}$, up to the polylogarithmic factor $(\log n)^{c_4}$ in the
upper bound. This proves \eqref{eq:rate}.
\end{proof}

\section{Concluding remarks}

\begin{enumerate}
\item \emph{Interpretation of the exponent.} The minimax rate can be written as
$n^{-2\alpha/(2\alpha+2p_{\mathrm{eff}}d+d)}$ with $2p_{\mathrm{eff}}d=3$: the problem
behaves like a mildly ill--posed problem whose total smoothing order is $3=2\times(d-1)
\text{ (angular)}+1\times1\text{ (radial)}$ spread over the $d$ dimensions. Equivalently,
the variance of a frequency--$L$ cut--off inversion grows like $L^{d+3}$, not
$L^{d+4}$ (isotropic degree $2$) nor $L^{d+2}$.

\item \emph{Comparison with boundary--constrained classes.} For the classes
$\{f\in H^\alpha: f|_{\partial D}\equiv1,\ \partial_\nu^jf=0,\ j\le\alpha-1\}$
studied in
\href{https://export.arxiv.org/pdf/1809.08818v3.pdf}{arXiv:1809.08818},
\href{https://arxiv.org/html/2409.03417}{arXiv:2409.03417},
\href{https://google.iopscience.iop.org/article/10.1088/1361-6420/ab7d2a}{Inverse Problems 36 (2020)}
and \href{https://arxiv.org/html/2601.09007}{arXiv:2601.09007}, the published recovery
rate for $f$ is $N^{-2(\alpha-1)/(2\alpha+d+2)}$, obtained from the minimax prediction
rate for $u$ and the (crude) $2$--ill--posed stability in $H^2$. Our results show that
the unconstrained problem has minimax rate $n^{-2\alpha/(2\alpha+d+3)}$; the two
statements are compatible (constrained classes are smaller, and their published rates
are upper bounds), and it is an interesting open problem to determine the exact minimax
rate in the constrained case.

\item \emph{The role of the ball.} The explicit solvability of $K\Lambda=\mathrm{Id}$
and the single critical point of $u_f$ near the origin are used essentially. On general
smooth domains the same transport mechanism operates along the gradient flow of $u_f$
between boundary and critical set, and we expect the same exponent whenever the critical
set is a nondegenerate minimum; we do not pursue this here.

\item \emph{Computational aspects.} The estimator \eqref{eq:estimator} requires only
(i) the spectral projection $P_L$, computable to accuracy $n^{-c}$ in time
$n^{\kappa}$, $\kappa<2$, by the wavelet discretisation schemes of
\href{https://arxiv.org/html/2601.09007}{arXiv:2601.09007}, and (ii) the one--dimensional
radial integrations defining $K$, which are $O(L^2)$ per angular mode. Hence the
minimax rate is attained by a polynomial--time algorithm.
\end{enumerate}

\begin{thebibliography}{9}

\bibitem{GT} D. Gilbarg, N. S. Trudinger, \emph{Elliptic Partial Differential Equations
of Second Order}, 2nd ed., Springer, 1983.

\bibitem{Triebel} H. Triebel, \emph{Theory of Function Spaces III}, Birkh\"auser, 1992.

\bibitem{Watson} G. N. Watson, \emph{A Treatise on the Theory of Bessel Functions},
2nd ed., Cambridge University Press, 1944.

\bibitem{Olver} F. W. J. Olver, \emph{Asymptotics and Special Functions}, AK Peters, 1997.

\bibitem{Tsybakov} A. B. Tsybakov, \emph{Introduction to Nonparametric Estimation},
Springer, 2009.

\bibitem{NvdG} R. Nickl, S. van de Geer, \emph{Convergence Rates for Penalized Least
Squares Estimators in PDE Constrained Regression Problems},
\href{https://export.arxiv.org/pdf/1809.08818v3.pdf}{arXiv:1809.08818}.

\bibitem{MAP24} \emph{Convergence Rates for the MAP Estimator in PDE-Regression Models
with Random Design}, \href{https://arxiv.org/html/2409.03417}{arXiv:2409.03417}.

\bibitem{GP20} \emph{Consistency of Bayesian inference with Gaussian process priors in
an elliptic inverse problem}, Inverse Problems 36 (2020),
\href{https://google.iopscience.iop.org/article/10.1088/1361-6420/ab7d2a}{iopscience.iop.org/article/10.1088/1361-6420/ab7d2a}.

\bibitem{PT26} \emph{Global polynomial-time estimation in statistical nonlinear inverse
problems via generalized stability},
\href{https://arxiv.org/html/2601.09007}{arXiv:2601.09007}.

\end{thebibliography}

\end{document}
```

\medskip
\noindent\textbf{Summary of the answer.} For the class $\F(\alpha,r)$ (a H\"older ball
around the constant conductivity $\mathbf 1$, \emph{without} boundary constraints) and
the white-noise model $Y_n=u_f+n^{-1/2}\W$, the minimax $L^2$ mean-squared error over
$D\subset\R^d$ satisfies, for small fixed $r$,
\[
\inf_{\hat f_n}\sup_{f_0\in\F(\alpha,r)}\E_{f_0}\norm{\hat f_n-f_0}_{L^2(D)}^2
\ \asymp\ r^{\frac{2(d+3)}{2\alpha+d+3}}\; n^{-\frac{2\alpha}{2\alpha+d+3}}
\quad\text{(up to log factors)},
\]
so the minimax rate is $\mathbf{n^{-2\alpha/(2\alpha+d+3)}}$ (i.e.\ $L^2$-error
$n^{-\alpha/(2\alpha+d+3)}$). The exponent $d+3$ reflects the anisotropy of the
linearised Darcy map at $f\equiv1$: degree-$2$ ill-posedness in the $d-1$ angular
directions but only degree-$1$ along the radial transport
$x\cdot\nabla h+dh=\divg(hx)$, whose exact inverse $Kh(r,\theta)=r^{-d}\int_0^r
h(s,\theta)s^{d-1}\ddd s$ (with $K\Lambda=\mathrm{Id}$ and $-dK\Delta$ a left inverse of
$DG_{\mathbf 1}$) drives both the upper bound (estimator $-dK\Delta P_LY_n$) and the
lower bound (Assouad over $\asymp L^{d-1}$ angular wave packets, each with inverse
attenuation $L^4$).